\documentclass{amsart}
\usepackage{amsmath, amssymb, amsthm}
\usepackage[margin=1in]{geometry}
\usepackage{enumitem}

\title{A Taste of Topology --- Chapter 2 Exercises}
\author{}
\date{}

\begin{document}
\maketitle

\subsection*{1}
An ant walks on the floor, ceiling, and walls of a cubical room. What metric is natural for the ant's view of its world? What metric would a spider consider natural? If the ant wants to walk from a point $p$ to a point $q$, how could it determine the shortest path?

\subsection*{2}
Why is the sum metric on $\mathbb{R}^2$ called the Manhattan metric and the taxicab metric?

\subsection*{3}
What is the set of points in $\mathbb{R}^3$ at distance exactly $1/2$ from the unit circle $S^1$ in the plane,
\[
T = \{p \in \mathbb{R}^3 : \exists\, q \in S^1 \text{ and } d(p,q) = 1/2 \text{ and for all } q' \in S^1 \text{ we have } d(p,q) \le d(p,q')\}?
\]

\subsection*{4}
Write out a proof that the discrete metric on a set $M$ is actually a metric.

\subsection*{5}
For $p, q \in S^1$, the unit circle in the plane, let
\[
d_a(p,q) = \min\{ |\angle(p) - \angle(q)|,\ 2\pi - |\angle(p) - \angle(q)| \}
\]
where $\angle(z) \in [0, 2\pi)$ refers to the angle that $z$ makes with the positive $x$-axis. Use your geometric talent to prove that $d_a$ is a metric on $S^1$.

\subsection*{6}
For $p, q \in [0, \pi/2]$ let
\[
d_s(p,q) = \sin|p-q|.
\]
Use your calculus talent to decide whether $d_s$ is a metric.

\subsection*{7}
Prove that every convergent sequence $(p_n)$ in a metric space $M$ is bounded, i.e., that for some $r>0$, some $q \in M$, and all $n \in \mathbb{N}$, we have $p_n \in M_r q$.

\subsection*{8}
Consider a sequence $(x_n)$ in the metric space $\mathbb{R}$.
\begin{enumerate}[label=(\alph*)]
\item If $(x_n)$ converges in $\mathbb{R}$ prove that the sequence of absolute values $(|x_n|)$ converges in $\mathbb{R}$.
\item State the converse.
\item Prove or disprove it.
\end{enumerate}

\subsection*{9}
A sequence $(x_n)$ in $\mathbb{R}$ increases if $n<m$ implies $x_n \le x_m$. It strictly increases if $n<m$ implies $x_n < x_m$. It decreases or strictly decreases if $n<m$ always implies $x_n \ge x_m$ or always implies $x_n > x_m$. A sequence is \textbf{monotone} if it increases or it decreases. Prove that every sequence in $\mathbb{R}$ which is monotone and bounded converges in $\mathbb{R}$.\footnote{This is nicely expressed by Pierre Teilhard de Chardin, ``Tout ce qui monte converge,'' in a different context.}

\subsection*{10}
Prove that the least upper bound property is equivalent to the ``monotone sequence property'' that every bounded monotone sequence converges.

\subsection*{11}
Let $(x_n)$ be a sequence in $\mathbb{R}$.
\begin{enumerate}[label=(\alph*)]
\item[*(a)] Prove that $(x_n)$ has a monotone subsequence.
\item How can you deduce that every bounded sequence in $\mathbb{R}$ has a convergent subsequence?
\item Infer that you have a second proof of the Bolzano--Weierstrass Theorem in $\mathbb{R}$.
\item What about the Heine--Borel Theorem?
\end{enumerate}

\subsection*{12}
Let $(p_n)$ be a sequence and $f : \mathbb{N} \to \mathbb{N}$ be a bijection. The sequence $(q_k)_{k \in \mathbb{N}}$ with $q_k = p_{f(k)}$ is a rearrangement of $(p_n)$.
\begin{enumerate}[label=(\alph*)]
\item Are limits of a sequence unaffected by rearrangement?
\item What if $f$ is an injection?
\item A surjection?
\end{enumerate}

\subsection*{13}
Assume that $f : M \to N$ is a function from one metric space to another which satisfies the following condition: If a sequence $(p_n)$ in $M$ converges then the sequence $(f(p_n))$ in $N$ converges. Prove that $f$ is continuous. [This result improves Theorem 4.]

\subsection*{14}
The simplest type of mapping from one metric space to another is an \textbf{isometry}. It is a bijection $f : M \to N$ that preserves distance in the sense that for all $p,q \in M$ we have
\[
d_N(fp, fq) = d_M(p,q).
\]
If there exists an isometry from $M$ to $N$ then $M$ and $N$ are said to be \textbf{isometric}, $M \equiv N$. You might have two copies of a unit equilateral triangle, one centered at the origin and one centered elsewhere. They are isometric. Isometric metric spaces are indistinguishable as metric spaces.
\begin{enumerate}[label=(\alph*)]
\item Prove that every isometry is continuous.
\item Prove that every isometry is a homeomorphism.
\item Prove that $[0,1]$ is not isometric to $[0,2]$.
\end{enumerate}

\subsection*{15}
Prove that isometry is an equivalence relation: If $M$ is isometric to $N$, show that $N$ is isometric to $M$; show that each $M$ is isometric to itself (what mapping of $M$ to $M$ is an isometry?); if $M$ is isometric to $N$ and $N$ is isometric to $P$, show that $M$ is isometric to $P$.

\subsection*{16}
Is the perimeter of a square isometric to the circle? Homeomorphic? Explain.

\subsection*{17}
Which capital letters of the Roman alphabet are homeomorphic? Are any isometric? Explain.

\subsection*{18}
Is $\mathbb{R}$ homeomorphic to $\mathbb{Q}$? Explain.

\subsection*{19}
Is $\mathbb{Q}$ homeomorphic to $\mathbb{N}$? Explain.

\subsection*{20}
What function (given by a formula) is a homeomorphism from $(-1,1)$ to $\mathbb{R}$? Is every open interval homeomorphic to $(0,1)$? Why or why not?

\subsection*{21}
Is the plane minus four points on the $x$-axis homeomorphic to the plane minus four points in an arbitrary configuration?

\subsection*{22}
If every closed and bounded subset of a metric space $M$ is compact, does it follow that $M$ is complete? (Proof or counterexample.)

\subsection*{23}
$(0,1)$ is an open subset of $\mathbb{R}$ but not of $\mathbb{R}^2$, when we think of $\mathbb{R}$ as the $x$-axis in $\mathbb{R}^2$. Prove this.

\subsection*{24}
For which intervals $[a,b]$ in $\mathbb{R}$ is the intersection $[a,b] \cap \mathbb{Q}$ a clopen subset of the metric space $\mathbb{Q}$?

\subsection*{25}
Prove directly from the definition of closed set that every singleton subset of a metric space $M$ is a closed subset of $M$. Why does this imply that every finite set of points is also a closed set?

\subsection*{26}
Prove that a set $U \subset M$ is open if and only if none of its points are limits of its complement.

\subsection*{27}
If $S, T \subset M$, a metric space, and $S \subset T$, prove that
\begin{enumerate}[label=(\alph*)]
\item $\overline{S} \subset \overline{T}$.
\item $\mathrm{int}(S) \subset \mathrm{int}(T)$.
\end{enumerate}

\subsection*{28}
A map $f : M \to N$ is \textbf{open} if for each open set $U \subset M$, the image set $f(U)$ is open in $N$.
\begin{enumerate}[label=(\alph*)]
\item If $f$ is open, is it continuous?
\item If $f$ is a homeomorphism, is it open?
\item If $f$ is an open, continuous bijection, is it a homeomorphism?
\item If $f : \mathbb{R} \to \mathbb{R}$ is a continuous surjection, must it be open?
\item If $f : \mathbb{R} \to \mathbb{R}$ is a continuous, open surjection, must it be a homeomorphism?
\item What happens in (e) if $\mathbb{R}$ is replaced by the unit circle $S^1$?
\end{enumerate}

\subsection*{29}
Let $\mathcal{T}$ be the collection of open subsets of a metric space $M$, and $\mathcal{K}$ the collection of closed subsets. Show that there is a bijection from $\mathcal{T}$ onto $\mathcal{K}$.

\subsection*{30}
Consider a two-point set $M = \{a,b\}$ whose topology consists of the two sets $M$ and the empty set. Why does this topology not arise from a metric on $M$?

\subsection*{31}
Prove the following.
\begin{enumerate}[label=(\alph*)]
\item If $U$ is an open subset of $\mathbb{R}$ then it consists of countably many disjoint intervals $U = \bigsqcup U_i$. (Unbounded intervals $(-\infty,b)$, $(a,\infty)$, and $(-\infty,\infty)$ are permitted.)
\item Prove that these intervals $U_i$ are uniquely determined by $U$. In other words, there is only one way to express $U$ as a disjoint union of open intervals.
\item If $U, V \subset \mathbb{R}$ are both open, so $U = \bigsqcup U_i$ and $V = \bigsqcup V_j$ where $U_i$ and $V_j$ are open intervals, show that $U$ and $V$ are homeomorphic if and only if there are equally many $U_i$ and $V_j$.
\end{enumerate}

\subsection*{32}
Show that every subset of $\mathbb{N}$ is clopen. What does this tell you about every function $f : \mathbb{N} \to M$, where $M$ is a metric space?

\subsection*{33}
\begin{enumerate}[label=(\alph*)]
\item Find a metric space in which the boundary of $M_r p$ is not equal to the sphere of radius $r$ at $p$, $\partial(M_r p) \ne \{x \in M : d(x,p) = r\}$.
\item Need the boundary be contained in the sphere?
\end{enumerate}

\subsection*{34}
Use the Inheritance Principle to prove Corollary 15.

\subsection*{35}
Prove that $S$ clusters at $p$ if and only if for each $r>0$ there is a point $q \in M_r p \cap S$, such that $q \ne p$.

\subsection*{36}
Construct a set with exactly three cluster points.

\subsection*{37}
Construct a function $f : \mathbb{R} \to \mathbb{R}$ that is continuous only at points of $\mathbb{Z}$.

\subsection*{38}
Let $X, Y$ be metric spaces with metrics $d_X, d_Y$, and let $M = X \times Y$ be their Cartesian product. Prove that the three natural metrics $d_E$, $d_{\max}$, and $d_{\mathrm{sum}}$ on $M$ are actually metrics. [Hint: Cauchy--Schwarz.]

\subsection*{39}
\begin{enumerate}[label=(\alph*)]
\item Prove that every convergent sequence is bounded. That is, if $(p_n)$ converges in the metric space $M$, prove that there is some neighborhood $M_r q$ containing the set $\{p_n : n \in \mathbb{N}\}$.
\item Is the same true for a Cauchy sequence in an incomplete metric space?
\end{enumerate}

\subsection*{40}
Let $M$ be a metric space with metric $d$. Prove that the following are equivalent.
\begin{enumerate}[label=(\alph*)]
\item $M$ is homeomorphic to $M$ equipped with the discrete metric.
\item Every function $f : M \to M$ is continuous.
\item Every bijection $g : M \to M$ is a homeomorphism.
\item $M$ has no cluster points.
\item Every subset of $M$ is clopen.
\item Every compact subset of $M$ is finite.
\end{enumerate}

\subsection*{41}
Let $\|\cdot\|$ be any norm on $\mathbb{R}^m$ and let $B = \{x \in \mathbb{R}^m : \|x\| \le 1\}$. Prove that $B$ is compact. [Hint: It suffices to show that $B$ is closed and bounded with respect to the Euclidean metric.]

\subsection*{42}
What is wrong with the following ``proof'' of Theorem 28? ``Let $((a_n,b_n))$ be any sequence in $A \times B$ where $A$ and $B$ are compact. Compactness implies the existence of subsequences $(a_{n_k})$ and $(b_{n_k})$ converging to $a \in A$ and $b \in B$ as $k \to \infty$. Therefore $((a_{n_k}, b_{n_k}))$ is a subsequence of $((a_n,b_n))$ that converges to a limit in $A \times B$, proving that $A \times B$ is compact.''

\subsection*{43}
Assume that the Cartesian product of two nonempty sets $A \subset M$ and $B \subset N$ is compact in $M \times N$. Prove that $A$ and $B$ are compact.

\subsection*{44}
Consider a function $f : M \to \mathbb{R}$. Its graph is the set
\[
\{(p,y) \in M \times \mathbb{R} : y = fp\}.
\]
\begin{enumerate}[label=(\alph*)]
\item Prove that if $f$ is continuous then its graph is closed (as a subset of $M \times \mathbb{R}$).
\item Prove that if $f$ is continuous and $M$ is compact then its graph is compact.
\item Prove that if the graph of $f$ is compact then $f$ is continuous.
\item What if the graph is merely closed? Give an example of a discontinuous function $f : \mathbb{R} \to \mathbb{R}$ whose graph is closed.
\end{enumerate}

\subsection*{45}
Draw a Cantor set $C$ on the circle and consider the set $A$ of all chords between points of $C$.
\begin{enumerate}[label=(\alph*)]
\item Prove that $A$ is compact.
\item[*(b)] Is $A$ convex?
\end{enumerate}

\subsection*{46}
Assume that $A, B$ are compact, disjoint, nonempty subsets of $M$. Prove that there are $a_0 \in A$ and $b_0 \in B$ such that for all $a \in A$ and $b \in B$ we have
\[
d(a_0,b_0) \le d(a,b).
\]
[The points $a_0, b_0$ are closest together.]

\subsection*{47}
Suppose that $A, B \subset \mathbb{R}^2$.
\begin{enumerate}[label=(\alph*)]
\item If $A$ and $B$ are homeomorphic, are their complements homeomorphic?
\item[*(b)] What if, in addition, $A$ and $B$ are compact?
\item[***(c)] What if $A$ and $B$ are compact and connected?
\end{enumerate}

\subsection*{48}
Prove that there is an embedding of the line as a closed subset of the plane, and there is an embedding of the line as a bounded subset of the plane, but there is no embedding of the line as a closed and bounded subset of the plane.

\subsection*{*49}
*Construct a subset $A \subset \mathbb{R}$ and a continuous bijection $f : A \to A$ that is not a homeomorphism. [Hint: By Theorem 36, $A$ must be noncompact.]

\subsection*{**50}
**Construct nonhomeomorphic connected, closed subsets $A, B \subset \mathbb{R}^2$ for which there exist continuous bijections $f : A \to B$ and $g : B \to A$. [Hint: By Theorem 36, $A$ and $B$ must be noncompact.]

\subsection*{***51}
***Do there exist nonhomeomorphic closed sets $A, B \subset \mathbb{R}$ for which there exist continuous bijections $f : A \to B$ and $g : B \to A$?

\subsection*{52}
Let $(A_n)$ be a nested decreasing sequence of nonempty closed sets in the metric space $M$.
\begin{enumerate}[label=(\alph*)]
\item If $M$ is complete and $\mathrm{diam}\, A_n \to 0$ as $n \to \infty$, show that $\bigcap A_n$ is exactly one point.
\item To what assertions do the sets $[n,\infty)$ provide counterexamples?
\end{enumerate}

\subsection*{53}
Suppose that $(K_n)$ is a nested sequence of compact nonempty sets, $K_1 \supset K_2 \supset \cdots$, and $K = \bigcap K_n$. If for some $\mu > 0$, $\mathrm{diam}\, K_n \ge \mu$ for all $n$, is it true that $\mathrm{diam}\, K \ge \mu$?

\subsection*{54}
If $f : A \to B$ and $g : C \to B$ such that $A \subset C$ and for each $a \in A$ we have $f(a) = g(a)$ then we also say that $g$ extends $f$. Assume that $f : S \to \mathbb{R}$ is a uniformly continuous function defined on a subset $S$ of a metric space $M$.
\begin{enumerate}[label=(\alph*)]
\item Prove that $f$ extends to a uniformly continuous function $\overline{f} : \overline{S} \to \mathbb{R}$.
\item Prove that $\overline{f}$ is the unique continuous extension of $f$ to a function defined on $\overline{S}$.
\item Prove the same things when $\mathbb{R}$ is replaced with a complete metric space $N$.
\end{enumerate}

\subsection*{55}
The distance from a point $p$ in a metric space $M$ to a nonempty subset $S \subset M$ is defined to be $\mathrm{dist}(p,S) = \inf\{d(p,s) : s \in S\}$.
\begin{enumerate}[label=(\alph*)]
\item Show that $p$ is a limit of $S$ if and only if $\mathrm{dist}(p,S) = 0$.
\item Show that $p \mapsto \mathrm{dist}(p,S)$ is a uniformly continuous function of $p \in M$.
\end{enumerate}

\subsection*{56}
Prove that the 2-sphere is not homeomorphic to the plane.

\subsection*{57}
If $S$ is connected, is the interior of $S$ connected? Prove this or give a counterexample.

\subsection*{58}
Theorem 49 states that the closure of a connected set is connected.
\begin{enumerate}[label=(\alph*)]
\item Is the closure of a disconnected set disconnected?
\item What about the interior of a disconnected set?
\end{enumerate}

\subsection*{*59}
*Prove that every countable metric space (not empty and not a singleton) is disconnected. [Astonishingly, there exists a countable topological space which is connected. Its topology does not arise from a metric.]

\subsection*{60}
\begin{enumerate}[label=(\alph*)]
\item Prove that a continuous function $f : M \to \mathbb{R}$, all of whose values are integers, is constant provided that $M$ is connected.
\item What if all the values are irrational?
\end{enumerate}

\subsection*{61}
Prove that the (double) cone $\{(x,y,z) \in \mathbb{R}^3 : x^2+y^2 = z^2\}$ is path-connected.

\subsection*{62}
Prove that the annulus $A = \{z \in \mathbb{R}^2 : r \le |z| \le R\}$ is connected.

\subsection*{63}
A subset $E$ of $\mathbb{R}^m$ is \textbf{starlike} if it contains a point $p_0$ (called a \textbf{center} for $E$) such that for each $q \in E$, the segment between $p_0$ and $q$ lies in $E$.
\begin{enumerate}[label=(\alph*)]
\item If $E$ is convex and nonempty prove that it is starlike.
\item Why is the converse false?
\item Is every starlike set connected?
\item Is every connected set starlike? Why or why not?
\end{enumerate}

\subsection*{*64}
*Suppose that $E \subset \mathbb{R}^m$ is open, bounded, and starlike, and $p_0$ is a center for $E$.
\begin{enumerate}[label=(\alph*)]
\item Is it true or false that all points $p_1$ in a small enough neighborhood of $p_0$ are also centers for $E$?
\item Is the set of centers convex?
\item Is it closed as a subset of $E$?
\item Can it consist of a single point?
\end{enumerate}

\subsection*{65}
Suppose that $A, B \subset \mathbb{R}^2$ are convex, closed, and have nonempty interiors.
\begin{enumerate}[label=(\alph*)]
\item Prove that $A, B$ are the closure of their interiors.
\item If $A, B$ are compact, prove that they are homeomorphic.
\end{enumerate}
[Hint: Draw a picture.]

\subsection*{66}
\begin{enumerate}[label=(\alph*)]
\item Prove that every connected open subset of $\mathbb{R}^m$ is path-connected.
\item Is the same true for open connected subsets of the circle?
\item What about connected nonopen subsets of the circle?
\end{enumerate}

\subsection*{67}
List the convex subsets of $\mathbb{R}$ up to homeomorphism. How many are there and how many are compact?

\subsection*{68}
List the closed convex sets in $\mathbb{R}^2$ up to homeomorphism. There are nine. How many are compact?

\subsection*{69}
Generalize Exercises 65 and 68 to $\mathbb{R}^3$; to $\mathbb{R}^m$.

\subsection*{70}
Prove that $(a,b)$ and $[a,b)$ are not homeomorphic metric spaces.

\subsection*{71}
Let $M$ and $N$ be nonempty metric spaces.
\begin{enumerate}[label=(\alph*)]
\item If $M$ and $N$ are connected prove that $M \times N$ is connected.
\item What about the converse?
\item Answer the questions again for path-connectedness.
\end{enumerate}

\subsection*{72}
Let $H$ be the hyperbola $\{(x,y) \in \mathbb{R}^2 : xy = 1 \text{ and } x,y > 0\}$ and let $X$ be the $x$-axis.
\begin{enumerate}[label=(\alph*)]
\item Is the set $S = X \cup H$ connected?
\item What if we replace $H$ with the graph $G$ of any continuous positive function $f : \mathbb{R} \to (0,\infty)$; is $X \cup G$ connected?
\item What if $f$ is everywhere positive but discontinuous at just one point.
\end{enumerate}

\subsection*{73}
Is the disc minus a countable set of points connected? Path-connected? What about the sphere or the torus instead of the disc?

\subsection*{74}
Let $S = \mathbb{R}^2 \setminus \mathbb{Q}^2$. (Points $(x,y) \in S$ have at least one irrational coordinate.) Is $S$ connected? Path-connected? Prove or disprove.

\subsection*{75}
An \textbf{arc} is a path with no self-intersection. Define the concept of arc-connectedness and prove that a metric space is path-connected if and only if it is arc-connected.

\subsection*{76}
\begin{enumerate}[label=(\alph*)]
\item The intersection of connected sets need not be connected. Give an example.
\item Suppose that $S_1, S_2, S_3, \ldots$ is a sequence of connected, closed subsets of the plane and $S_1 \supset S_2 \supset \cdots$. Is $S = \bigcap S_n$ connected? Give a proof or counterexample.
\item[*(c)] Does the answer change if the sets are compact?
\item What is the situation for a nested decreasing sequence of compact path-connected sets?
\end{enumerate}

\subsection*{77}
If a metric space $M$ is the union of path-connected sets $S_\alpha$, all of which have the nonempty path-connected set $K$ in common, is $M$ path-connected?

\subsection*{78}
$(p_1, \ldots, p_n)$ is an $\epsilon$-chain in a metric space $M$ if for each $i$ we have $p_i \in M$ and $d(p_i, p_{i+1}) < \epsilon$. The metric space is \textbf{chain-connected} if for each $\epsilon > 0$ and each pair of points $p,q \in M$ there is an $\epsilon$-chain from $p$ to $q$.
\begin{enumerate}[label=(\alph*)]
\item Show that every connected metric space is chain-connected.
\item Show that if $M$ is compact and chain-connected then it is connected.
\item Is $\mathbb{R} \setminus \mathbb{Z}$ chain-connected?
\item If $M$ is complete and chain-connected, is it connected?
\end{enumerate}

\subsection*{79}
Prove that if $M$ is nonempty, compact, locally path-connected, and connected then it is path-connected. (See Exercise 143, below.)

\subsection*{80}
The \textbf{Hawaiian earring} is the union of circles of radius $1/n$ and center $\pm 1/n$ on the $x$-axis, for $n \in \mathbb{N}$. See Figure 27 on page 58.
\begin{enumerate}[label=(\alph*)]
\item Is it connected?
\item Path-connected?
\item Is it homeomorphic to the one-sided Hawaiian earring?
\end{enumerate}

\subsection*{*81}
*The topologist's sine curve is the set
\[
\{(x,y) : x=0 \text{ and } |y|\le 1\} \cup \{(x,y) : 0 < x \le 1 \text{ and } y = \sin(1/x)\}.
\]
See Figure 43. The \textbf{topologist's sine circle} is shown in Figure 58. (It is the union of a circular arc and the topologist's sine curve.) Prove that it is path-connected but not locally path-connected. ($M$ is \textbf{locally path-connected} if for each $p \in M$ and each neighborhood $U$ of $p$ there is a path-connected subneighborhood $V$ of $p$.)

\subsection*{82}
The graph of $f : M \to \mathbb{R}$ is the set $\{(x,y) \in M \times \mathbb{R} : y = fx\}$.
\begin{enumerate}[label=(\alph*)]
\item If $M$ is connected and $f$ is continuous, prove that the graph of $f$ is connected.
\item Give an example to show that the converse is false.
\item If $M$ is path-connected and $f$ is continuous, show that the graph is path-connected.
\item What about the converse?
\end{enumerate}

\subsection*{83}
The open cylinder is $(0,1) \times S^1$. The punctured plane is $\mathbb{R}^2 \setminus \{0\}$.
\begin{enumerate}[label=(\alph*)]
\item Prove that the open cylinder is homeomorphic to the punctured plane.
\item Prove that the open cylinder, the double cone, and the plane are not homeomorphic.
\end{enumerate}

\subsection*{84}
Is the closed strip $\{(x,y) \in \mathbb{R}^2 : 0 \le x \le 1\}$ homeomorphic to the closed half-plane $\{(x,y) \in \mathbb{R}^2 : x \ge 0\}$? Prove or disprove.

\subsection*{85}
Suppose that $M$ is compact and that $\mathcal{U}$ is an open covering of $M$ which is ``redundant'' in the sense that each $p \in M$ is contained in at least two members of $\mathcal{U}$. Show that $\mathcal{U}$ reduces to a finite subcovering with the same property.

\subsection*{86}
Suppose that every open covering of $M$ has a positive Lebesgue number. Give an example of such an $M$ that is not compact.

\subsection*{87}
Give a direct proof that $[a,b]$ is covering compact. [Hint: Let $\mathcal{U}$ be an open covering of $[a,b]$ and consider the set
\[
C = \{x \in [a,b] : \text{finitely many members of } \mathcal{U} \text{ cover } [a,x]\}.
\]
Use the least upper bound principle to show that $b \in C$.]

\subsection*{88}
Give a direct proof that a closed subset $A$ of a covering compact set $K$ is covering compact. [Hint: If $\mathcal{U}$ is an open covering of $A$, adjoin the set $W = M \setminus A$ to $\mathcal{U}$. Is $\mathcal{W} = \mathcal{U} \cup \{W\}$ an open covering of $K$? If so, so what?]

\subsection*{89}
Give a proof of Theorem 36 using open coverings. That is, assume $A$ is a covering compact subset of $M$ and $f : M \to N$ is continuous. Prove directly that $fA$ is covering compact. [Hint: What is the criterion for continuity in terms of preimages?]

\subsection*{90}
Suppose that $f : M \to N$ is a continuous bijection and $M$ is covering compact. Prove directly that $f$ is a homeomorphism.

\subsection*{91}
Suppose that $M$ is covering compact and that $f : M \to N$ is continuous. Use the Lebesgue number lemma to prove that $f$ is uniformly continuous. [Hint: Consider the covering of $N$ by $\epsilon/2$-neighborhoods $\{N_{\epsilon/2}(q) : q \in N\}$ and its preimage in $M$, $\{f^{\mathrm{pre}}(N_{\epsilon/2}(q)) : q \in N\}$.]

\subsection*{92}
Give a direct proof that the nested decreasing intersection of nonempty covering compact sets is nonempty. [Hint: If $A_1 \supset A_2 \supset \cdots$ are covering compact, consider the open sets $U_n = A_n^c$. If $\bigcap A_n = \emptyset$, what does $\{U_n\}$ cover?]

\subsection*{93}
Generalize Exercise 92 as follows. Suppose that $M$ is covering compact and $\mathcal{C}$ is a collection of closed subsets of $M$ such that every intersection of finitely many members of $\mathcal{C}$ is nonempty. (Such a collection $\mathcal{C}$ is said to have the \textbf{finite intersection property}.) Prove that the grand intersection $\bigcap_{C \in \mathcal{C}} C$ is nonempty. [Hint: Consider the collection of open sets $\mathcal{U} = \{C^c : C \in \mathcal{C}\}$.]

\subsection*{94}
If every collection of closed subsets of $M$ which has the finite intersection property also has a nonempty grand intersection, prove that $M$ is covering compact. [Hint: Given an open covering $\mathcal{U} = \{U_\alpha\}$, consider the collection of closed sets $\mathcal{C} = \{U_\alpha^c\}$.]

\subsection*{95}
Let $S$ be a subset of a metric space $M$. With respect to the definitions on page 92 prove the following.
\begin{enumerate}[label=(\alph*)]
\item The closure of $S$ is the intersection of all closed subsets of $M$ that contain $S$.
\item The interior of $S$ is the union of all open subsets of $M$ that are contained in $S$.
\item The boundary of $S$ is a closed set.
\item Why does (a) imply the closure of $S$ equals $\lim S$?
\item If $S$ is clopen, what is $\partial S$?
\item Give an example of $S \subset \mathbb{R}$ such that $\partial(\partial S) \ne \emptyset$, and infer that ``the boundary of the boundary $\partial \circ \partial$ is not always zero.''
\end{enumerate}

\subsection*{96}
If $A \subset B \subset C$, $A$ is dense in $B$, and $B$ is dense in $C$, prove that $A$ is dense in $C$.

\subsection*{97}
Is the set of dyadic rationals (the denominators are powers of 2) dense in $\mathbb{Q}$? In $\mathbb{R}$? Does one answer imply the other? (Recall that $A$ is dense in $B$ if $A \subset B$ and $\overline{A} \supset B$.)

\subsection*{98}
Show that $S \subset M$ is somewhere dense in $M$ if and only if $\mathrm{int}(\overline{S}) \ne \emptyset$. Equivalently, $S$ is nowhere dense in $M$ if and only if its closure has empty interior.

\subsection*{99}
Let $M, N$ be nonempty metric spaces and $P = M \times N$.
\begin{enumerate}[label=(\alph*)]
\item If $M, N$ are perfect prove that $P$ is perfect.
\item If $M, N$ are totally disconnected prove that $P$ is totally disconnected.
\item What about the converses?
\item Infer that the Cartesian product of Cantor spaces is a Cantor space. (We already know that the Cartesian product of compacts is compact.)
\item Why does this imply that $C \times C = \{(x,y) \in \mathbb{R}^2 : x \in C \text{ and } y \in C\}$ is homeomorphic to $C$, $C$ being the standard Cantor set?
\end{enumerate}

\subsection*{100}
Prove that every Cantor piece is a Cantor space. (Recall that $M$ is a Cantor space if it is compact, nonempty, totally disconnected and perfect, and that $A \subset M$ is a Cantor piece if it is nonempty and clopen.)

\subsection*{*101}
*Let $\Sigma$ be the set of all infinite sequences of zeroes and ones. For example, $(1001110000111\ldots) \in \Sigma$. Define the metric
\[
d(a,b) = \sum \frac{|a_n - b_n|}{2^n}
\]
where $a=(a_n)$ and $b=(b_n)$ are points in $\Sigma$.
\begin{enumerate}[label=(\alph*)]
\item Prove that $\Sigma$ is compact.
\item Prove that $\Sigma$ is homeomorphic to the Cantor set.
\end{enumerate}

\subsection*{102}
Prove that no Peano curve is one-to-one. (Recall that a Peano curve is a continuous map $f : [0,1] \to \mathbb{R}^2$ whose image has a nonempty interior.)

\subsection*{103}
Prove that there is a continuous surjection $\mathbb{R} \to \mathbb{R}^2$. What about $\mathbb{R}^m$?

\subsection*{104}
Find two nonhomeomorphic compact subsets of $\mathbb{R}$ whose complements are homeomorphic.

\subsection*{105}
As on page 115, consider the subsets of $\mathbb{R}$,
\[
A = \{0\} \cup [1,2] \cup \{3\} \quad \text{and} \quad B = \{0\} \cup \{1\} \cup [2,3].
\]
\begin{enumerate}[label=(\alph*)]
\item Why is there no ambient homeomorphism of $\mathbb{R}$ to itself that carries $A$ onto $B$?
\item Thinking of $\mathbb{R}$ as the $x$-axis, is there an ambient homeomorphism of $\mathbb{R}^2$ to itself that carries $A$ onto $B$?
\end{enumerate}

\subsection*{106}
Prove that the completion of a metric space is unique in the following natural sense: A completion of a metric space $M$ is a complete metric space $X$ containing $M$ as a metric subspace such that $M$ is dense in $X$. That is, every point of $X$ is a limit of $M$.
\begin{enumerate}[label=(\alph*)]
\item Prove that $M$ is dense in the completion $\widehat{M}$ constructed in the proof of Theorem 80.
\item If $X$ and $X'$ are two completions of $M$ prove that there is an isometry $i : X \to X'$ such that $i(p) = p$ for all $p \in M$.
\item Prove that $i$ is the unique such isometry.
\item Infer that $\widehat{M}$ is unique.
\end{enumerate}

\subsection*{107}
If $M$ is a metric subspace of a complete metric space $S$ prove that $\overline{M}$ is a completion of $M$.

\subsection*{*108}
*Consider the identity map $\mathrm{id} : C_{\max} \to C_{\mathrm{int}}$ where $C_{\max}$ is the metric space $C([0,1],\mathbb{R})$ of continuous real-valued functions defined on $[0,1]$, equipped with the max-metric $d_{\max}(f,g) = \max|f(x)-g(x)|$, and $C_{\mathrm{int}}$ is $C([0,1],\mathbb{R})$ equipped with the integral metric,
\[
d_{\mathrm{int}}(f,g) = \int_0^1 |f(x)-g(x)|\,dx.
\]
Show that $\mathrm{id}$ is a continuous linear bijection (an isomorphism) but its inverse is not continuous.

\subsection*{*109}
*A metric on $M$ is an \textbf{ultrametric} if for all $x,y,z \in M$ we have
\[
d(x,z) \le \max\{d(x,y), d(y,z)\}.
\]
(Intuitively this means that the trip from $x$ to $z$ cannot be broken into shorter legs by making a stopover at some $y$.)
\begin{enumerate}[label=(\alph*)]
\item Show that the ultrametric property implies the triangle inequality.
\item In an ultrametric space show that ``all triangles are isosceles.''
\item Show that a metric space with an ultrametric is totally disconnected.
\item Define a metric on the set $\Sigma$ of strings of zeroes and ones in Exercise 101 as
\[
d_*(a,b) =
\begin{cases}
\dfrac{1}{2^n} & \text{if } n \text{ is the smallest index for which } a_n \ne b_n \\[6pt]
0 & \text{if } a = b.
\end{cases}
\]
Show that $d_*$ is an ultrametric and prove that the identity map is a homeomorphism $(\Sigma, d) \to (\Sigma, d_*)$.
\end{enumerate}

\subsection*{*110}
*$\mathbb{Q}$ inherits the Euclidean metric from $\mathbb{R}$ but it also carries a very different metric, the $p$-adic metric. Given a prime number $p$ and an integer $n$, the $p$-adic norm of $n$ is
\[
|n|_p = \frac{1}{p^k}
\]
where $p^k$ is the largest power of $p$ that divides $n$. (The norm of $0$ is by definition $0$.) The more factors of $p$, the smaller the $p$-norm. Similarly, if $x = a/b$ is a fraction, we factor $x$ as
\[
x = p^k \cdot \frac{r}{s}
\]
where $p$ divides neither $r$ nor $s$, and we set
\[
|x|_p = \frac{1}{p^k}.
\]
The $p$-adic metric on $\mathbb{Q}$ is
\[
d_p(x,y) = |x-y|_p.
\]
\begin{enumerate}[label=(\alph*)]
\item Prove that $d_p$ is a metric with respect to which $\mathbb{Q}$ is perfect --- every point is a cluster point.
\item Prove that $d_p$ is an ultrametric.
\item Let $\mathbb{Q}_p$ be the metric space completion of $\mathbb{Q}$ with respect to the metric $d_p$, and observe that the extension of $d_p$ to $\mathbb{Q}_p$ remains an ultrametric. Infer from Exercise 109 that $\mathbb{Q}_p$ is totally disconnected.
\item Prove that $\mathbb{Q}_p$ is locally compact, in the sense that every point has small compact neighborhoods.
\item Infer that $\mathbb{Q}_p$ is covered by neighborhoods homeomorphic to the Cantor set. See Gouv\^ea's book, \emph{p-adic Numbers}.
\end{enumerate}

\subsection*{111}
Let $M=[0,1]$ and let $\mathcal{M}_1$ be its division into two intervals $[0,1/2]$ and $[1/2,1]$. Let $\mathcal{M}_2$ be its division into four intervals $[0,1/4]$, $[1/4,1/2]$, $[1/2,3/4]$, and $[3/4,1]$. Continuing these bisections generates natural divisions of $[0,1]$. The pieces are intervals. We label them with words using the letters 0 and 1 as follows: 0 means ``left'' and 1 means ``right,'' so the four intervals in $\mathcal{M}_2$ are labeled as 00, 01, 10, and 11 respectively.
\begin{enumerate}[label=(\alph*)]
\item Verify that all endpoints of the intervals (except 0 and 1) have two addresses. For instance,
\[
\bigcap_k \left[\frac{2^{k-1}-1}{2^k}, \frac{1}{2}\right] = \left\{\frac{1}{2}\right\} = \bigcap_k \left[\frac{1}{2}, \frac{2^{k-1}+1}{2^k}\right].
\]
\item Verify that the points $0, 1$, and all nonendpoints have unique addresses.
\end{enumerate}

\subsection*{*112}
*Prove that $\#C = \#\mathbb{R}$. [Hint: According to the Schroeder--Bernstein Theorem from Chapter 1 it suffices to find injections $C \to \mathbb{R}$ and $\mathbb{R} \to C$. The inclusion $C \subset \mathbb{R}$ is an injection $C \to \mathbb{R}$. Each $t \in [0,1)$ has a unique base-2 expansion $\tau(t)$ that does not terminate in an infinite string of ones. Replacing each 1 by 2 converts $\tau(t)$ to $\omega(t)$, an infinite address in the symbols 0 and 2. It does not terminate in an infinite string of twos. Set $h(t) = \sum_{i=1}^\infty \omega_i/3^i$ and verify that $h : [0,1) \to C$ is an injection. Since there is an injection $\mathbb{R} \to [0,1)$, conclude that there is an injection $\mathbb{R} \to C$, and hence that $\#C = \#\mathbb{R}$.]

\textbf{Remark.} The Continuum Hypothesis states that if $S$ is any uncountable subset of $\mathbb{R}$ then $S$ and $\mathbb{R}$ have equal cardinality. The preceding coding shows that $C$ is not only uncountable (as is implied by Theorem 56) but actually has the same cardinality as $\mathbb{R}$. That is, $C$ is not a counterexample to the Continuum Hypothesis. The same is true of all uncountable closed subsets of $\mathbb{R}$. See Exercise 151.

\subsection*{113}
Let $M$ be the standard Cantor set $C$. In the notation of Section 8, $C^n$ is the collection of $2^n$ Cantor intervals of length $1/3^n$ that nest down to $C$ as $n \to \infty$. Verify that setting $\mathcal{C}_k = C \cap C^k$ gives divisions of $C$ into disjoint clopen pieces.

\subsection*{*114}
*
\begin{enumerate}[label=(\alph*)]
\item Prove directly that there is a continuous surjection of the middle-thirds Cantor set $C$ onto the closed interval $[0,1]$. [Hint: Each $x \in C$ has a base-3 expansion $(x_n)$, all of whose entries are zeroes and twos. (For example, $2/3 = (2\overline{0})_{\text{base }3}$ and $1/3 = (0\overline{2})_{\text{base }3}$.) Write $y = (y_n)$ by replacing the twos in $(x_n)$ by ones and interpreting the answer base 2. Show that the map $x \mapsto y$ works.]
\item Compare this surjection to the one constructed from the bisection divisions in Exercise 113.
\end{enumerate}

\subsection*{115}
Rotate the unit circle $S^1$ by a fixed angle $\alpha$, say $R : S^1 \to S^1$. (In polar coordinates, the transformation $R$ sends $(1,\theta)$ to $(1,\theta+\alpha)$.)
\begin{enumerate}[label=(\alph*)]
\item If $\alpha/\pi$ is rational, show that each orbit of $R$ is a finite set.
\item[*(b)] If $\alpha/\pi$ is irrational, show that each orbit is infinite and has closure equal to $S^1$.
\end{enumerate}

\subsection*{116}
A metric space $M$ with metric $d$ can always be remetrized so the metric becomes bounded. Simply define the bounded metric
\[
\rho(p,q) = \frac{d(p,q)}{1+d(p,q)}.
\]
\begin{enumerate}[label=(\alph*)]
\item Prove that $\rho$ is a metric. Why is it obviously bounded?
\item Prove that the identity map $M \to M$ is a homeomorphism from $M$ with the $d$-metric to $M$ with the $\rho$-metric.
\item Infer that boundedness of $M$ is not a topological property.
\item Find homeomorphic metric spaces, one bounded and the other not.
\end{enumerate}

\subsection*{117}
Fold a piece of paper in half.
\begin{enumerate}[label=(\alph*)]
\item Is this a continuous transformation of one rectangle into another?
\item Is it injective?
\item Draw an open set in the target rectangle, and find its preimage in the original rectangle. Is it open?
\item What if the open set meets the crease?
\end{enumerate}
The \textbf{baker's transformation} is a similar mapping. A rectangle of dough is stretched to twice its length and then folded back on itself. Is the transformation continuous? A formula for the baker's transformation in one variable is $f(x) = 1-|1-2x|$. The $n$th iterate of $f$ is $f^n = f \circ f \circ \cdots \circ f$, $n$ times. The \textbf{orbit} of a point $x$ is
\[
\{x, f(x), f^2(x), \ldots, f^n(x), \ldots\}.
\]
[For clearer but more awkward notation one can write $f^{\circ n}$ instead of $f^n$. This distinguishes composition $f \circ f$ from multiplication $f \cdot f$.]
\begin{enumerate}[resume,label=(\alph*)]
\item If $x$ is rational prove that the orbit of $x$ is a finite set.
\item If $x$ is irrational what is the orbit?
\end{enumerate}

\subsection*{*118}
*The implications of compactness are frequently equivalent to it. Prove
\begin{enumerate}[label=(\alph*)]
\item If every continuous function $f : M \to \mathbb{R}$ is bounded then $M$ is compact.
\item If every continuous bounded function $f : M \to \mathbb{R}$ achieves a maximum or minimum then $M$ is compact.
\item If every continuous function $f : M \to \mathbb{R}$ has compact range $fM$ then $M$ is compact.
\item If every nested decreasing sequence of nonempty closed subsets of $M$ has nonempty intersection then $M$ is compact.
\end{enumerate}
Together with Theorems 63 and 65, (a)--(d) give seven equivalent definitions of compactness. [Hint: Reason contrapositively. If $M$ is not compact then it contains a sequence $(p_n)$ that has no convergent subsequence. It is fair to assume that the points $p_n$ are distinct. Find radii $r_n>0$ such that the neighborhoods $M_{r_n}(p_n)$ are disjoint and no sequence $q_n \in M_{r_n}(p_n)$ has a convergent subsequence. Using the metric define a function $f_n : M_{r_n}(p_n) \to \mathbb{R}$ with a spike at $p_n$, such as
\[
f_n(x) = \frac{r_n - d(x,p_n)}{a_n + d(x,p_n)}
\]
where $a_n>0$. Set $f(x) = f_n(x)$ if $x \in M_{r_n}(p_n)$, and $f(x)=0$ if $x$ belongs to no $M_{r_n}(p_n)$. Show that $f$ is continuous. With the right choice of $a_n$ show that $f$ is unbounded. With a different choice of $a_n$, it is bounded but achieves no maximum, and so on.]

\subsection*{119}
Let $M$ be a metric space of diameter $\le 2$. The cone for $M$ is the set
\[
C = C(M) = \{p_0\} \cup M \times (0,1]
\]
with the cone metric
\[
\rho((p,s),(q,t)) = |s-t| + \min\{s,t\}\, d(p,q), \qquad \rho((p,s),p_0) = s, \qquad \rho(p_0,p_0) = 0.
\]
The point $p_0$ is the vertex of the cone. Prove that $\rho$ is a metric on $C$. [If $M$ is the unit circle, think of it in the plane $z=1$ in $\mathbb{R}^3$ centered at the point $(0,0,1)$. Its cone is the 45-degree cone with vertex the origin.]

\subsection*{120}
Recall that if for each embedding of $M$, $h : M \to N$, $hM$ is closed in $N$ then $M$ is said to be \textbf{absolutely closed}. If each $hM$ is bounded then $M$ is \textbf{absolutely bounded}. Theorem 41 implies that compact sets are absolutely closed and absolutely bounded. Prove:
\begin{enumerate}[label=(\alph*)]
\item If $M$ is absolutely bounded then $M$ is compact.
\item[*(b)] If $M$ is absolutely closed then $M$ is compact.
\end{enumerate}
Thus these are two more conditions equivalent to compactness. [Hint: From Exercise 118(a), if $M$ is noncompact there is a continuous function $f : M \to \mathbb{R}$ that is unbounded. For Exercise 120(a), show that $F(x) = (x,f(x))$ embeds $M$ onto a nonbounded subset of $M \times \mathbb{R}$. For 120(b), justify the additional assumption that the metric on $M$ is bounded by 2. Then use Exercise 118(b) to show that if $M$ is noncompact then there is a continuous function $g : M \to (0,1]$ such that for some nonclustering sequence $(p_n)$, we have $g(p_n) \to 0$ as $n \to \infty$. Finally, show that $G(x) = (x,gx)$ embeds $M$ onto a nonclosed subset $S$ of the cone $C(M)$ discussed in Exercise 119. $S$ will be nonclosed because it limits at $p_0$ but does not contain it.]

\subsection*{121}
\begin{enumerate}[label=(\alph*)]
\item Prove that every function defined on a discrete metric space is uniformly continuous.
\item Infer that it is false to assert that if every continuous function $f : M \to \mathbb{R}$ is uniformly continuous then $M$ is compact.
\item Prove, however, that if $M$ is a metric subspace of a compact metric space $K$ and every continuous function $f : M \to \mathbb{R}$ is uniformly continuous then $M$ is compact.
\end{enumerate}

\subsection*{122}
Recall that $p$ is a cluster point of $S$ if each $M_r p$ contains infinitely many points of $S$. The set of cluster points of $S$ is denoted as $S'$. Prove:
\begin{enumerate}[label=(\alph*)]
\item If $S \subset T$ then $S' \subset T'$.
\item $(S \cup T)' = S' \cup T'$.
\item $S' = (\overline{S})'$.
\item $S'$ is closed in $M$; that is, $S'' \subset S'$ where $S'' = (S')'$.
\item Calculate $\mathbb{N}'$, $\mathbb{Q}'$, $\mathbb{R}'$, $(\mathbb{R}\setminus\mathbb{Q})'$, and $\mathbb{Q}''$.
\item Let $T$ be the set of points $\{1/n : n \in \mathbb{N}\}$. Calculate $T'$ and $T''$.
\item Give an example showing that $S''$ can be a proper subset of $S'$.
\end{enumerate}

\subsection*{123}
Recall that $p$ is a condensation point of $S$ if each $M_r p$ contains uncountably many points of $S$. The set of condensation points of $S$ is denoted as $S^*$. Prove:
\begin{enumerate}[label=(\alph*)]
\item If $S \subset T$ then $S^* \subset T^*$.
\item $(S \cup T)^* = S^* \cup T^*$.
\item $S^* \subset \overline{S}^*$ where $\overline{S}^* = (\overline{S})^*$.
\item $S^*$ is closed in $M$; that is, $S^{*\prime} \subset S^*$ where $S^{*\prime} = (S^*)'$.
\item $S^{**} \subset S^*$ where $S^{**} = (S^*)^*$.
\item Calculate $\mathbb{N}^*$, $\mathbb{Q}^*$, $\mathbb{R}^*$, and $(\mathbb{R}\setminus\mathbb{Q})^*$.
\item Give an example showing that $S^*$ can be a proper subset of $(\overline{S})^*$. Thus, (c) is not in general an equality.
\item[**(h)] Give an example that $S^{**}$ can be a proper subset of $S^*$. Thus, (e) is not in general an equality. [Hint: Consider the set $M$ of all functions $f : [a,b] \to [0,1]$, continuous or not, and let the metric on $M$ be the sup metric, $d(f,g) = \sup\{|f(x)-g(x)| : x \in [a,b]\}$. Consider the set $S$ of all ``$\delta$-functions with rational values.'']
\item[**(i)] Give examples that show in general that $S^*$ neither contains nor is contained in $S'^*$ where $S'^* = (S')^*$. [Hint: $\delta$-functions with values $1/n$, $n \in \mathbb{N}$.]
\end{enumerate}

\subsection*{124}
Recall that $p$ is an interior point of $S \subset M$ if some $M_r p$ is contained in $S$. The set of interior points of $S$ is the interior of $S$ and is denoted $\mathrm{int}\, S$. For all subsets $S,T$ of the metric space $M$ prove:
\begin{enumerate}[label=(\alph*)]
\item $\mathrm{int}\, S = S \setminus \partial S$.
\item $\mathrm{int}\, S = (\overline{S^c})^c$.
\item $\mathrm{int}(\mathrm{int}\, S) = \mathrm{int}\, S$.
\item $\mathrm{int}(S \cap T) = \mathrm{int}(S \cap \mathrm{int}\, T)$.
\item What are the dual equations for the closure?
\item Prove that $\mathrm{int}(S \cup T) \supset \mathrm{int}\, S \cup \mathrm{int}\, T$. Show by example that the inclusion can be strict, i.e., not an equality.
\end{enumerate}

\subsection*{125}
A point $p$ is a boundary point of a set $S \subset M$ if every neighborhood $M_r p$ contains points of both $S$ and $S^c$. The boundary of $S$ is denoted $\partial S$. For all subsets $S, T$ of a metric space $M$ prove:
\begin{enumerate}[label=(\alph*)]
\item $S$ is clopen if and only if $\partial S = \emptyset$.
\item $\partial S = \partial S^c$.
\item $\partial \partial S \subset \partial S$.
\item $\partial \partial \partial S = \partial \partial S$.
\item $\partial(S \cup T) \subset \partial S \cup \partial T$.
\item Give an example in which (c) is a strict inclusion, $\partial \partial S \ne \partial S$.
\item What about (e)?
\end{enumerate}

\subsection*{*126}
*Suppose that $E$ is an uncountable subset of $\mathbb{R}$. Prove that there exists a point $p \in \mathbb{R}$ at which $E$ condenses. [Hint: Use decimal expansions. Why must there be an interval $[n,n+1)$ containing uncountably many points of $E$? Why must it contain a decimal subinterval with the same property? (A decimal subinterval $[a,b]$ has endpoints $a = n+k/10$, $b = n+(k+1)/10$ for some digit $k$, $0 \le k \le 9$.) Do you see lurking the decimal expansion of a condensation point?] Generalize to $\mathbb{R}^2$ and to $\mathbb{R}^m$.

\subsection*{127}
The metric space $M$ is \textbf{separable} if it contains a countable dense subset. [Note the confusion of language: ``Separable'' has nothing to do with ``separation.'']
\begin{enumerate}[label=(\alph*)]
\item Prove that $\mathbb{R}^m$ is separable.
\item Prove that every compact metric space is separable.
\end{enumerate}

\subsection*{128}
*
\begin{enumerate}[label=(\alph*)]
\item Prove that every metric subspace of a separable metric space is separable, and deduce that every metric subspace of $\mathbb{R}^m$ or of a compact metric space is separable.
\item Is the property of being separable topological?
\item Is the continuous image of a separable metric space separable?
\end{enumerate}

\subsection*{129}
Think up a nonseparable metric space.

\subsection*{130}
Let $\mathcal{B}$ denote the collection of all $\epsilon$-neighborhoods in $\mathbb{R}^m$ whose radius $\epsilon$ is rational and whose center has all coordinates rational.
\begin{enumerate}[label=(\alph*)]
\item Prove that $\mathcal{B}$ is countable.
\item Prove that every open subset of $\mathbb{R}^m$ can be expressed as the countable union of members of $\mathcal{B}$.
\end{enumerate}
(The union need not be disjoint, but it is at most a countable union because there are only countably many members of $\mathcal{B}$. A collection such as $\mathcal{B}$ is called a \textbf{countable base} for the topology of $\mathbb{R}^m$.)

\subsection*{131}
\begin{enumerate}[label=(\alph*)]
\item Prove that every separable metric space has a countable base for its topology, and conversely that every metric space with a countable base for its topology is separable.
\item Infer that every compact metric space has a countable base for its topology.
\end{enumerate}

\subsection*{*132}
*Referring to Exercise 123, assume now that $M$ is separable, $S \subset M$, and, as before $S'$ is the set of cluster points of $S$ while $S^*$ is the set of condensation points of $S$. Prove:
\begin{enumerate}[label=(\alph*)]
\item $S^* \subset (S')^* = (\overline{S})^*$.
\item $S^{**} = S^{*\prime} = S^*$.
\item Why is (a) not in general an equality?
\end{enumerate}
[Hints: For (a) write $S \subset (S \setminus S') \cup S'$ and $\overline{S} = (S \setminus S') \cup S'$, show that $(S \setminus S')^* = \emptyset$, and use Exercise 123(a). For (b), Exercise 123(d) implies that $S^{**} \subset S^{*\prime} \subset S^*$. To prove that $S^* \subset S^{**}$, write $S \subset (S \setminus S^*) \cup S^*$ and show that $(S \setminus S^*)^* = \emptyset$.]

\subsection*{*133}
*Prove that
\begin{enumerate}[label=(\alph*)]
\item An uncountable subset of $\mathbb{R}$ clusters at some point of $\mathbb{R}$.
\item An uncountable subset of $\mathbb{R}$ clusters at some point of itself.
\item An uncountable subset of $\mathbb{R}$ condenses at uncountably many points of itself.
\item What about $\mathbb{R}^m$ instead of $\mathbb{R}$?
\item What about any compact metric space?
\item What about any separable metric space?
\end{enumerate}
[Hint: Review Exercise 126.]

\subsection*{*134}
*Prove that $\widehat{\mathbb{Q}}$, the Cauchy sequences in $\mathbb{Q}$ modulo the equivalence relation of being co-Cauchy, is a field with respect to the natural arithmetic operations defined on page 122, and that $\mathbb{Q}$ is naturally a subfield of $\widehat{\mathbb{Q}}$.

\subsection*{135}
Prove that the order on $\widehat{\mathbb{Q}}$ defined on page 122 is a bona fide order which agrees with the standard order on $\mathbb{Q}$.

\subsection*{*136}
*Let $M$ be the square $[0,1]^2$, and let $aa, ba, bb, ab$ label its four quadrants --- upper right, upper left, lower left, and lower right.
\begin{enumerate}[label=(\alph*)]
\item Define nested bisections of the square using this pattern repeatedly, and let $\tau_k$ be a curve composed of line segments that visit the $k$th-order quadrants systematically. Let $\tau = \lim_k \tau_k$ be the resulting Peano curve \`a la the Cantor Surjection Theorem.
\item Compare $\tau$ to the Peano curve $f : I \to I^2$ directly constructed on pages 271--274 of the second edition of Munkres' book \emph{Topology}.
\end{enumerate}

\subsection*{*137}
*Let $P$ be a closed perfect subset of a separable complete metric space $M$. Prove that each point of $P$ is a condensation point of $P$. In symbols, $P = P' \Rightarrow P = P^*$.

\subsection*{**138}
**Given a Cantor space $M \subset \mathbb{R}^2$, given a line segment $[p,q] \subset \mathbb{R}^2$ with $p,q \notin M$, and given an $\epsilon>0$, prove that there exists a path $A$ in the $\epsilon$-neighborhood of $[p,q]$ that joins $p$ to $q$ and is disjoint from $M$. [Hint: Think of $A$ as a bisector of $M$. From this bisection fact a dyadic disc partition of $M$ can be constructed, which leads to the proof that $M$ is tame.]

\subsection*{139}
To prove that Antoine's Necklace $A$ is a Cantor set, you need to show that $A$ is compact, perfect, nonempty, and totally disconnected.
\begin{enumerate}[label=(\alph*)]
\item Do so. [Hint: What is the diameter of any connected component of $A^n$, and what does that imply about $A$?]
\item[**(b)] If, in the Antoine construction two linked solid tori are placed \emph{very cleverly} inside each larger solid torus, show that the intersection $A = \bigcap A^n$ is a Cantor set.
\end{enumerate}

\subsection*{*140}
*Consider the Hilbert cube
\[
H = \left\{ (x_1,x_2,\ldots) \in [0,1]^\infty : \text{for each } n \in \mathbb{N} \text{ we have } |x_n| \le 1/2^n \right\}.
\]
Prove that $H$ is compact with respect to the metric
\[
d(x,y) = \sup_n |x_n - y_n|
\]
where $x=(x_n)$, $y=(y_n)$. [Hint: Sequences of sequences.]

\textbf{Remark.} Although compact, $H$ is infinite-dimensional and is homeomorphic to no subset of $\mathbb{R}^m$.

\subsection*{141}
Prove that the Hilbert cube is perfect and homeomorphic to its Cartesian square, $H \cong H \times H$.

\subsection*{***142}
***Assume that $M$ is compact, nonempty, perfect, and homeomorphic to its Cartesian square, $M \cong M \times M$. Must $M$ be homeomorphic to the Cantor set, the Hilbert cube, or some combination of them?

\subsection*{143}
A \textbf{Peano space} is a metric space $M$ that is the continuous image of the unit interval: There is a continuous surjection $\tau : [0,1] \to M$. Theorem 72 states the amazing fact that the 2-disc is a Peano space. Prove that every Peano space is
\begin{enumerate}[label=(\alph*)]
\item compact,
\item nonempty,
\item path-connected,
\item[*(d)] and locally path-connected, in the sense that for each $p \in M$ and each neighborhood $U$ of $p$ there is a smaller neighborhood $V$ of $p$ such that any two points of $V$ can be joined by a path in $U$.
\end{enumerate}

\subsection*{*144}
*The converse to Exercise 143 is the \textbf{Hahn--Mazurkiewicz Theorem}. Assume that a metric space $M$ is compact, nonempty, path-connected, and locally path-connected. Use the Cantor Surjection Theorem 70 to show that $M$ is a Peano space. [The key is to make uniformly short paths to fill in the gaps of $[0,1] \setminus C$.]

\subsection*{145}
One of the famous theorems in plane topology is the \textbf{Jordan Curve Theorem}. A Jordan curve $J$ is a homeomorph of the unit circle in the plane. (Equivalently it is $f([a,b])$ where $f : [a,b] \to \mathbb{R}^2$ is continuous, $f(a)=f(b)$, and for no other pair of distinct $s,t \in [a,b]$ does $f(s)$ equal $f(t)$. It is also called a \textbf{simple closed curve}.) The Jordan Curve Theorem asserts that $\mathbb{R}^2 \setminus J$ consists of two disjoint, connected open sets, its inside and its outside, and every path between them must meet $J$. Prove the Jordan Curve Theorem for the circle, the square, the triangle, and --- if you have courage --- every simple closed polygon.

\subsection*{146}
The \textbf{utility problem} gives three houses 1, 2, 3 in the plane and the three utilities, Gas, Water, and Electricity. You are supposed to connect each house to the three utilities without crossing utility lines. (The houses and utilities are disjoint.)
\begin{enumerate}[label=(\alph*)]
\item Use the Jordan curve theorem to show that there is no solution to the utility problem in the plane.
\item[*(b)] Show also that the utility problem cannot be solved on the 2-sphere $S^2$.
\item[*(c)] Show that the utility problem can be solved on the surface of the torus.
\item[*(d)] What about the surface of the Klein bottle?
\item[***(e)] Given utilities $U_1,\ldots,U_m$ and houses $H_1,\ldots,H_n$ located on a surface with $g$ handles, find necessary and sufficient conditions on $m,n,g$ so that the utility problem can be solved.
\end{enumerate}

\subsection*{147}
Let $M$ be a metric space and let $\mathcal{K}$ denote the class of nonempty compact subsets of $M$. The $r$-neighborhood of $A \in \mathcal{K}$ is
\[
M_r A = \{x \in M : \exists\, a \in A \text{ and } d(x,a) < r\} = \bigcup_{a \in A} M_r a.
\]
For $A, B \in \mathcal{K}$ define
\[
D(A,B) = \inf\{r>0 : A \subset M_r B \text{ and } B \subset M_r A\}.
\]
\begin{enumerate}[label=(\alph*)]
\item Show that $D$ is a metric on $\mathcal{K}$. (It is called the \textbf{Hausdorff metric} and $\mathcal{K}$ is called the \textbf{hyperspace} of $M$.)
\item Denote by $\mathcal{F}$ the collection of finite nonempty subsets of $M$ and prove that $\mathcal{F}$ is dense in $\mathcal{K}$. That is, given $A \in \mathcal{K}$ and given $\epsilon>0$ show there exists $F \in \mathcal{F}$ such that $D(A,F) < \epsilon$.
\item[*(c)] If $M$ is compact prove that $\mathcal{K}$ is compact.
\item If $M$ is connected prove that $\mathcal{K}$ is connected.
\item[**(e)] If $M$ is path-connected is $\mathcal{K}$ path-connected?
\item Do homeomorphic metric spaces have homeomorphic hyperspaces?
\end{enumerate}
\textbf{Remark.} The converse to (f), $\mathcal{K}(M) \cong \mathcal{K}(N) \Rightarrow M \cong N$ is false. The hyperspace of every Peano space is the Hilbert cube. This is a difficult result but a good place to begin reading about hyperspaces is Sam Nadler's book \emph{Continuum Theory}.

\subsection*{**148}
**Start with a set $S \subset \mathbb{R}$ and successively take its closure, the complement of its closure, the closure of that, and so on: $S, \mathrm{cl}(S), (\mathrm{cl}(S))^c, \ldots$. Do the same to $S^c$. In total, how many distinct subsets of $\mathbb{R}$ can be produced this way? In particular decide whether each chain $S, \mathrm{cl}(S), \ldots$ consists of only finitely many sets. For example, if $S = \mathbb{Q}$ then we get $\mathbb{Q}, \mathbb{R}, \emptyset, \emptyset, \mathbb{R}, \mathbb{R}, \ldots$ and $\mathbb{Q}^c, \mathbb{R}, \emptyset, \emptyset, \mathbb{R}, \mathbb{R}, \ldots$ for a total of four sets.

\subsection*{**149}
**Consider the letter T.
\begin{enumerate}[label=(\alph*)]
\item Prove that there is no way to place uncountably many copies of the letter T disjointly in the plane. [Hint: First prove this when the unit square replaces the plane.]
\item Prove that there is no way to place uncountably many homeomorphic copies of the letter T disjointly in the plane.
\item For which other letters of the alphabet is this true?
\item Let $U$ be a set in $\mathbb{R}^3$ formed like an umbrella: It is a disc with a perpendicular segment attached to its center. Prove that uncountably many copies of $U$ cannot be placed disjointly in $\mathbb{R}^3$.
\item What if the perpendicular segment is attached to the boundary of the disc?
\end{enumerate}

\subsection*{**150}
**Let $M$ be a complete, separable metric space such as $\mathbb{R}^m$. Prove the \textbf{Cupcake Theorem}: Each closed set $K \subset M$ can be expressed uniquely as the disjoint union of a countable set and a perfect closed set,
\[
C \sqcup P = K.
\]

\subsection*{**151}
**Let $M$ be an uncountable compact metric space.
\begin{enumerate}[label=(\alph*)]
\item Prove that $M$ contains a homeomorphic copy of the standard Cantor set $C$. [Hint: Imitate the construction of the standard Cantor set.]
\item Infer that Cantor sets are ubiquitous. There is a continuous surjection $\sigma : C \to M$ and there is a continuous injection $i : C \to M$.
\item Infer that every uncountable closed set $S \subset \mathbb{R}$ has $\#S = \#\mathbb{R}$, and hence that the Continuum Hypothesis is valid for closed sets in $\mathbb{R}$. [Hint: Cupcake and Exercise 112.]
\item Is the same true if $M$ is separable, uncountable, and complete?
\end{enumerate}

\end{document}
